\chapter{Premises}
\label{ch:premises}

A design is a set of choices made subject to constraints, and most bad designs
for Cuba are bad because they quietly relax a constraint that does not relax.
This chapter fixes the constraint set. Each item is a finding of Parts I and II
rather than an assumption, each is stated in the form that matters for the
chapters ahead, and each is something the design must work \emph{with} rather
than around.

\section{The population is the binding scarcity}

Cuba has fewer than ten million residents and is losing more. Fertility has been
below replacement since 1978, so there is no cohort in the pipeline; more than a
quarter of the population is over sixty; and between one and two million people,
disproportionately of working and child-bearing age, left between 2021 and 2025.

The consequence is a reversal of the assumption most development planning makes.
Cuba is not a country with surplus labour looking for capital to employ it. It is
a country with a shrinking, ageing workforce, and every proposal in Part III has
to be read as competing for the same scarce people. A hotel that opens draws
staff from a hospital. A remote-work sector that grows draws engineers from the
state. This is why the book does not propose labour-intensive
import-substitution, why Chapter~\ref{ch:welfare} treats the pension arithmetic
as its hardest problem, and why Chapter~\ref{ch:talent} treats immigration and
return migration as fiscal policy rather than as sentiment.

It also sets the clock. No policy adopted in 2026 changes the working-age
population before the 2040s except through migration, and the emigration is
self-reinforcing: each departure removes a reason for the next person to stay
and adds a household abroad that can finance their leaving. The design's first
phase is not judged on growth. It is judged on whether the outflow slows.

\section{The state is insolvent and has no access to credit}

Cuba has been in default on its convertible-currency debt since 1986. The 2015
Paris Club arrangement, which forgave a large part of accumulated arrears and
rescheduled the rest, broke down within four years when Cuba stopped paying
under it. Russia wrote off ninety per cent of the Soviet-era debt in 2014.
Suppliers have unpaid invoices going back to 2009, when the state froze foreign
companies' accounts. Cuba is not a member of the International Monetary Fund,
the World Bank or the Inter-American Development Bank, and so has access neither
to their lending nor to the technical assistance and implicit certification
that other countries use to attract private capital.

Two things follow. Nothing in Part III can be financed by sovereign borrowing on
ordinary terms, because there are no such terms available; the financing has to
come from equity, from the diaspora, from concessional and bilateral sources,
and above all from cash flow. And the sequence in Chapter~\ref{ch:money} has to
put debt normalisation early, not because settling it releases money --- it does
not --- but because until it is settled every other transaction in the country
is priced as though the counterparty defaults, which it has.

\section{Nothing works until the electricity does}

This is the reason Chapter~\ref{ch:energy} precedes every sector chapter.

Electricity is not one sector among several in Cuba; it is the input to all of
them. Water is pumped. Food is refrigerated. Hospitals cannot run lists without
it. Tourism cannot be sold without it --- a hotel without power or water does
not merely disappoint a guest, it destroys the review that produces the next
guest. Remote work is impossible without it. Manufacturing is impossible without
it. And a grid that has collapsed nationally four times collapses more easily
each time.

The design premise is therefore blunt: any plan whose sector chapters would work
only if the power were reliable is a plan that has to deliver reliable power
first, and its first eighteen months belong to that and to very little else.

\section{The welfare floor does not come down}

This is stated in the preface and repeated here as a premise because it binds
everything after it.

The design retains a universal welfare state at Western European standard ---
health, education, pensions, disability, income support --- and the rest of Part
III is arranged around being able to afford it. The reasoning is in
Chapter~\ref{ch:welfare}. What belongs here is the cost, admitted in advance
rather than discovered in a later chapter: a European benefit structure requires
a tax burden well above the Latin American norm, which is a real competitive
disadvantage against every other country in the region bidding for the same
investment. This book resolves that trade-off in favour of the welfare state and
compensates on the investment side with certainty, speed and rule of law rather
than with low rates. Chapter~\ref{ch:capital} is written on that basis and must
not quietly assume the opposite.

The floor is also, in 2026, thinner than it has ever been and depreciating.
Retaining it is not a matter of leaving something alone. It requires spending
that does not currently exist, on staff who are currently leaving.

\section{The diaspora is an asset with a grievance}

Something over two million people of Cuban origin live abroad, most within
three hours' flight, holding capital, professional skills, business networks
and market access that no policy can otherwise buy. Their remittances are
already among the largest sources of foreign exchange in the country. They are
the single largest potential source of the investment Chapter~\ref{ch:capital}
needs, and they are the natural first market for
Chapter~\ref{ch:tourism}'s visitors and Chapter~\ref{ch:talent}'s clients.

They also hold claims and they hold memories, and the memories are specific
rather than diffuse: confiscated houses, the inventory taken at the door, the
years of agricultural labour that were the price of an exit application, the
children of Peter Pan, the mobs outside the door in 1980, and the relatives who
did not get out. A design that treats the diaspora as a chequebook will fail,
and a design that treats every claim as valid and every asset as returnable will
fail differently and worse. Chapter~\ref{ch:capital} proposes a bounded,
non-restitutionary settlement, and Chapter~\ref{ch:polity} proposes a
truth-telling mechanism with no prosecutorial power, precisely because these two
things have to be settled together and neither can be settled by generosity in
one direction only.

\section{The incumbents are in the room}

This is the premise most reform proposals for Cuba omit, and omitting it is what
makes them unserious.

Any transition from the present arrangement is designed and executed by, or at
minimum with the acquiescence of, three groups that currently hold power: the
Communist Party apparatus, the Revolutionary Armed Forces, and the military-
linked business conglomerate GAESA, which controls a large share of the
country's hard-currency earnings through tourism, retail and logistics --- how
large is not published, which is itself the point.

These are not abstractions to be legislated away. They hold the assets, the
guns, the foreign-currency accounts and the administrative capacity, and they
have watched what happened to the Soviet nomenklatura, to Ceau\c{s}escu, and to
Gaddafi. A design that offers them nothing is a design they will not permit, and
the historical record is unambiguous that they can prevent it. A design that
offers them everything is a design that produces an oligarchy with a flag, which
Chapter~\ref{ch:integrity} argues is the most likely failure mode and the one
that is effectively irreversible once it happens.

The book therefore treats the incumbents as a party to be bargained with
explicitly: guarantees of personal security, pensions and property in exchange
for the surrender of discretionary control, with the bargain written down,
sequenced, and made verifiable. Chapters~\ref{ch:polity},
\ref{ch:capital} and \ref{ch:integrity} each carry part of it.

\section{The United States is a variable, not a constant}

The embargo may be relaxed, tightened, or left where it is, and Cuban policy
does not control which. Helms--Burton converted it from an executive instrument
into a statute, so its removal requires an act of Congress rather than a
decision by a president, and the domestic American politics of that act have
been immovable for thirty years and are not obviously about to move.

The design premise is symmetric and it is the only defensible one: \emph{assume
nothing}. Nothing in Part III may be contingent on the embargo ending, because
it may not; and nothing in Part III may be built on the assumption that it will
persist, because a sudden normalisation is also possible and a country with no
plan for it will be overrun by the resulting inflow --- of capital, of visitors,
of claims --- exactly as it was overrun in miniature in 2015--16. Each chapter
that touches the American relationship therefore states what it does under both
branches.

\section{Race and region decide who is exposed}

Cuba's inequality has a colour and a compass bearing, and both were made by
policy.

The remittance economy flows overwhelmingly to households whose relatives left
in the 1960s, who were overwhelmingly white. The tourism sector's
foreign-currency jobs have been allocated in part by appearance. Meanwhile the
five eastern provinces --- the poorest, the most agricultural, the most
Afro-Cuban --- have the worst housing, the worst infrastructure and the longest
blackouts, and are the origin of most internal migration to Havana, which the
state has restricted by decree since 1997.

Every liberalising measure in Part III advantages the household with capital,
with a relative abroad, with a house in a tourist district, with a computer.
That is not an argument against the measures; it is an argument for knowing in
advance who is on the other side of them, and for the compensating instruments
--- the universal floor of Chapter~\ref{ch:welfare}, the regional allocation
rules in Chapters~\ref{ch:energy} and \ref{ch:tourism}, and the deliberately
broad-based structure of Chapter~\ref{ch:capital}'s privatisation --- to be
designed in from the start rather than added when the disparity becomes
visible. Cuba has already run the experiment of declaring the racial question
solved and then not measuring it for thirty years. The design's first
requirement in this domain is to publish the data.

\section{What this book does not assume}

Finally, the negative premises, stated so that a reader can hold the argument to
them.

\emph{No shock therapy.} The book does not propose simultaneous liberalisation
of all prices, rapid mass privatisation, or the withdrawal of the social state
in advance of the institutions that would replace it. The historical record of
that package is bad everywhere and would be worse here, for the demographic
reason in the first section: a country whose young have already left does not
survive a decade of transitional depression.

\emph{No restitution by expulsion.} No Cuban is removed from a dwelling to
return it to a prior owner. Chapter~\ref{ch:capital} settles claims in money and
paper, within a cap.

\emph{No assumption about Washington.} As above, in both directions.

\emph{No assumption of regime collapse, and none of regime permanence.} The
design is written so that its early phases are executable by the present
government if it chooses, and its later phases by whatever government follows.
That is deliberate. A plan that requires a particular political rupture before
its first measure can be taken is a plan that does nothing until the rupture,
and Cubans have been waiting for the rupture since 1991.

\emph{No claim of certainty.} Chapter~\ref{ch:sequence} closes with a list of
indicators by which the design can be judged to be failing. A proposal that
cannot be falsified is not a proposal.
